% This document was generated by an LLM.

\documentclass{essay}

\usepackage{siunitx}
\DeclareSIUnit{\var}{var}

\newcommand{\ee}{\mathrm{e}}
\newcommand{\jj}{\mathrm{j}}
\renewcommand{\Re}{\operatorname{Re}}
\renewcommand{\Im}{\operatorname{Im}}

\title{AC and Transient Analysis}
\subtitle{A Formula Reference for Linear Circuits --- Inductors, Capacitors,
Impedance, Charge, Flux, Voltage \& Current}
\author{}
\date{}

\begin{document}

\maketitle

This reference collects the constitutive relations of passive elements, the
governing equations of first- and second-order transients, and the
phasor/impedance description of sinusoidal steady state. Throughout, the
passive sign convention is used: current $i$ enters the terminal at which the
voltage $v$ is taken to be positive. The imaginary unit is written $\jj$ so as
not to clash with current, and $\omega = 2\pi f$ denotes angular frequency.

\section{Constitutive Relations of Passive Elements}

The three ideal two-terminal elements are defined by how voltage and current
relate at their terminals.

\subsection{Resistor}
A resistor obeys Ohm's law instantaneously; it stores no energy.
\[
  v(t) = R\,i(t), \qquad i(t) = \frac{v(t)}{R}, \qquad p(t)=v\,i =
  i^2 R = \frac{v^2}{R}.
\]

\subsection{Capacitor --- charge and current}
A capacitor stores energy in its electric field. Its defining variable is the
stored charge $q$, proportional to the terminal voltage through the capacitance
$C$:
\[
  q(t) = C\,v(t), \qquad
  i(t) = \frac{\dd q}{\dd t} = C\,\frac{\dd v}{\dd t}, \qquad
  v(t) = \frac{1}{C}\int_{-\infty}^{t} i(\tau)\,\dd\tau .
\]
With a known initial voltage $v(t_0)$,
\[
  v(t) = v(t_0) + \frac{1}{C}\int_{t_0}^{t} i(\tau)\,\dd\tau .
\]
The voltage across a capacitor cannot change instantaneously for a finite
current, so $v_C$ is a continuous state variable.

\subsection{Inductor --- flux and voltage}
An inductor stores energy in its magnetic field. Its defining variable is the
flux linkage $\lambda$, proportional to the current through the inductance $L$:
\[
  \lambda(t) = L\,i(t), \qquad
  v(t) = \frac{\dd \lambda}{\dd t} = L\,\frac{\dd i}{\dd t}, \qquad
  i(t) = \frac{1}{L}\int_{-\infty}^{t} v(\tau)\,\dd\tau .
\]
With a known initial current $i(t_0)$,
\[
  i(t) = i(t_0) + \frac{1}{L}\int_{t_0}^{t} v(\tau)\,\dd\tau .
\]
The current through an inductor cannot change instantaneously for a finite
voltage, so $i_L$ is a continuous state variable.

\subsection{Stored energy}
\[
  W_C(t) = \tfrac{1}{2}\,C\,v^2 = \frac{q^2}{2C}, \qquad
  W_L(t) = \tfrac{1}{2}\,L\,i^2 = \frac{\lambda^2}{2L}.
\]

\begin{center}
  \begin{tabular}{@{}lccc@{}}
    \hline
    & \textbf{Resistor} & \textbf{Capacitor} & \textbf{Inductor}\\
    \hline
    $v$--$i$ law      & $v=Ri$                 & $i=C\,\dd v/\dd t$
    & $v=L\,\dd i/\dd t$\\
    State variable    & ---                    & $v_C$ (charge
    $q=Cv$)    & $i_L$ (flux $\lambda=Li$)\\
    Stored energy     & $0$                    & $\tfrac12 C v^2$
    & $\tfrac12 L i^2$\\
    Cannot jump       & ---                    & voltage
    & current\\
    \hline
  \end{tabular}
\end{center}

\section{First-Order Transients}

A circuit with a single energy-storage element reduces, via Th\'evenin
reduction of the resistive part, to a first-order linear ODE. Let $x(t)$ be the
state variable ($v_C$ or $i_L$).

\subsection{General first-order form}
\[
  \frac{\dd x}{\dd t} + \frac{1}{\tau}\,x = \frac{1}{\tau}\,x_\infty,
  \qquad
  x(t) = x_\infty + \bigl[x(0^+) - x_\infty\bigr]\,\ee^{-t/\tau}.
\]
Here $x(0^+)$ is the value just after the switching instant, $x_\infty$ is the
final (steady-state) value, and $\tau$ is the time constant.

\subsection{Time constants}
\[
  \tau_{RC} = R_{\mathrm{th}}\,C, \qquad \tau_{RL} = \frac{L}{R_{\mathrm{th}}} .
\]
$R_{\mathrm{th}}$ is the Th\'evenin resistance seen by the storage element with
independent sources deactivated.

\subsection{RC circuit: charging and discharging}
For a capacitor charged through $R$ from a source $V_s$ (initially uncharged):
\[
  v_C(t) = V_s\!\left(1-\ee^{-t/RC}\right), \qquad
  i(t) = \frac{V_s}{R}\,\ee^{-t/RC}.
\]
For a capacitor discharging from $V_0$ through $R$:
\[
  v_C(t) = V_0\,\ee^{-t/RC}, \qquad
  i(t) = -\frac{V_0}{R}\,\ee^{-t/RC}.
\]

\subsection{RL circuit: energizing and de-energizing}
For an inductor energized through $R$ from $V_s$ (initially zero current):
\[
  i_L(t) = \frac{V_s}{R}\!\left(1-\ee^{-Rt/L}\right), \qquad
  v_L(t) = V_s\,\ee^{-Rt/L}.
\]
For an inductor de-energizing from $I_0$ through $R$:
\[
  i_L(t) = I_0\,\ee^{-Rt/L}, \qquad
  v_L(t) = -I_0 R\,\ee^{-Rt/L}.
\]

After $5\tau$ the response is within $0.7\%$ of its final value and is treated
as settled.

\section{Second-Order Transients (Series \& Parallel RLC)}

With both $L$ and $C$ present the natural response satisfies a second-order
ODE. Writing the source-free equation for the state variable $x$:
\[
  \frac{\dd^2 x}{\dd t^2} + 2\alpha\,\frac{\dd x}{\dd t} + \omega_0^2\,x = 0,
  \qquad
  s^2 + 2\alpha s + \omega_0^2 = 0,
  \qquad
  s_{1,2} = -\alpha \pm \sqrt{\alpha^2 - \omega_0^2}.
\]

\subsection{Defining parameters}
\[
  \omega_0 = \frac{1}{\sqrt{LC}} \;\;(\text{undamped natural freq.}),
  \qquad
  \zeta = \frac{\alpha}{\omega_0} \;\;(\text{damping ratio}),
  \qquad
  Q = \frac{1}{2\zeta}.
\]
\begin{center}
  \begin{tabular}{@{}lcc@{}}
    \hline
    \textbf{Quantity} & \textbf{Series RLC} & \textbf{Parallel RLC}\\
    \hline
    Neper frequency $\alpha$ & $\dfrac{R}{2L}$ & $\dfrac{1}{2RC}$\\
    Resonant freq. $\omega_0$ & $\dfrac{1}{\sqrt{LC}}$ &
    $\dfrac{1}{\sqrt{LC}}$\\
    Quality factor $Q$ & $\dfrac{1}{R}\sqrt{\dfrac{L}{C}}$ &
    $R\sqrt{\dfrac{C}{L}}$\\
    \hline
  \end{tabular}
\end{center}

\subsection{Damping cases}
\begin{itemize}
  \item \textbf{Overdamped} ($\alpha > \omega_0$):
    $x(t) = A_1\ee^{s_1 t} + A_2\ee^{s_2 t}$.
  \item \textbf{Critically damped} ($\alpha = \omega_0$):
    $x(t) = (A_1 + A_2 t)\,\ee^{-\alpha t}$.
  \item \textbf{Underdamped} ($\alpha < \omega_0$):
    $x(t) = \ee^{-\alpha t}\bigl(A_1\cos\omega_d t + A_2\sin\omega_d t\bigr)$,
    with damped frequency $\omega_d = \sqrt{\omega_0^2 - \alpha^2}$.
\end{itemize}

\section{Sinusoidal Steady State and Phasors}

A sinusoid $x(t) = X_m\cos(\omega t + \phi)$ is represented by the phasor
$\mathbf{X} = X_m\,\ee^{\jj\phi} = X_m\angle\phi$, recovered through
\[
  x(t) = \Re\!\left\{\mathbf{X}\,\ee^{\jj\omega t}\right\}.
\]
Time differentiation maps to multiplication by $\jj\omega$, and integration to
division by $\jj\omega$:
\[
  \frac{\dd}{\dd t} \;\longleftrightarrow\; \jj\omega, \qquad
  \int \dd t \;\longleftrightarrow\; \frac{1}{\jj\omega}.
\]

\subsection{Impedance and admittance}
Impedance $Z = V/I$ generalizes resistance to the complex frequency domain;
admittance is its reciprocal $Y = 1/Z = G + \jj B$.
\[
  Z_R = R, \qquad
  Z_L = \jj\omega L, \qquad
  Z_C = \frac{1}{\jj\omega C} = -\frac{\jj}{\omega C}.
\]
In rectangular and polar form,
\[
  Z = R + \jj X = |Z|\,\ee^{\jj\theta}, \qquad
  |Z| = \sqrt{R^2 + X^2}, \qquad
  \theta = \arctan\!\frac{X}{R},
\]
where $X$ is the reactance: $X_L = \omega L > 0$ for an inductor and
$X_C = -1/(\omega C) < 0$ for a capacitor.

\subsection{Series and parallel combination}
\[
  Z_{\text{series}} = \sum_k Z_k, \qquad
  \frac{1}{Z_{\text{parallel}}} = \sum_k \frac{1}{Z_k}, \qquad
  Z_{1\parallel 2} = \frac{Z_1 Z_2}{Z_1 + Z_2}.
\]

\subsection{Resonance}
Series and parallel RLC circuits resonate when the reactive parts cancel,
$\omega_0 L = 1/(\omega_0 C)$:
\[
  \omega_0 = \frac{1}{\sqrt{LC}}, \qquad
  f_0 = \frac{1}{2\pi\sqrt{LC}}, \qquad
  \text{bandwidth } \mathrm{BW} = \frac{\omega_0}{Q}.
\]

\section{AC Power}

For voltage and current phasors $\mathbf{V}=V_m\angle\theta_v$,
$\mathbf{I}=I_m\angle\theta_i$, define the phase difference
$\theta = \theta_v - \theta_i$ and RMS magnitudes
$V_{\mathrm{rms}}=V_m/\sqrt2$, $I_{\mathrm{rms}}=I_m/\sqrt2$.

\[
  \mathbf{S} = \mathbf{V}_{\mathrm{rms}}\,\mathbf{I}_{\mathrm{rms}}^{*}
  = P + \jj Q,
  \qquad
  |\mathbf{S}| = V_{\mathrm{rms}} I_{\mathrm{rms}}.
\]
\[
  P = V_{\mathrm{rms}} I_{\mathrm{rms}}\cos\theta \;\;(\text{real, \si{\watt}}),
  \qquad
  Q = V_{\mathrm{rms}} I_{\mathrm{rms}}\sin\theta
  \;\;(\text{reactive, \si{\var}}).
\]
The power factor is $\mathrm{pf} = \cos\theta = P/|\mathbf{S}|$ --- lagging for
inductive loads ($\theta>0$) and leading for capacitive loads ($\theta<0$).

\end{document}
